\documentclass[11pt]{article}
\renewcommand{\baselinestretch}{1.20} 
\usepackage[utf8]{inputenc}
\usepackage[english]{babel}
\usepackage{graphicx}
\usepackage{wrapfig}
\usepackage{subcaption}
\usepackage{geometry}
\usepackage{xcolor}
\usepackage{float}
\usepackage{mdframed}
\usepackage{fancyhdr}
\usepackage{lastpage}
\usepackage{booktabs}% http://ctan.org/pkg/booktabs
\newcommand{\tabitem}{~~\llap{\textbullet}~~}
\geometry{a4paper, total={170mm,237mm}, left=20mm, top=20mm}
\setlength{\abovecaptionskip}{15pt plus 3pt minus 2pt}

\pagestyle{fancy}
\fancyhf{}
\rfoot{Side \thepage \hspace{1pt} / \pageref{LastPage}}
\begin{document}
    
    % Title page
    \begin{titlepage}
    \centering
	\includegraphics[width=0.35\textwidth]{Projectdoc/Assets/Illustrationer/aau_logo_en.pdf}\par\vspace{1cm}
	{\scshape\Large Struktureret System Udvikling\par}
	\vspace{0.2cm}
	{\huge\bfseries Workshop 1\par}
	\vspace{0.2cm}
	{\scshape\Large ITC - B125\par}
	\vspace{2cm}
	{\Large\itshape 
    	Mikkel Steen Hansen\\
        Martin Boe\\
        Benjamin Bach Jensen\\
        Daniel Vestergaard Jensen\\
    \par}
	\vfill
	\vfill
\end{titlepage}
    
    % Table of contents
    \renewcommand{\baselinestretch}{0.8}
    \tableofcontents
    \renewcommand{\baselinestretch}{1.20}
    
    \section{Beskrivelse}
    \textbf{Mini projekt: Vejr station}\\
    Dagens workshop omhandler at videre design på sidste workshops vejr station for anvendelse i et parcel hus. \\
    Tanken er at forskellige sensorer registrerer en eller flere informationer der skal til for at beskrive
    vejret. Disse data samles op via en enhed der kan videregive information til en app på brugerens smart
    phone eller computer, evt. krydret med en præcis vejrudsigt i den kommende time(r) ved at udnytte services fra DMI
    eller anden vejr tjeneste.
    
    Gennem denne del har gruppen valgt at omskrive fokus fra tidligere Workshop. Dette er grundet at tidligere Workshop havde fokus på projektets login funktionalitet og de dertilhørende usecases. Til denne workshop er dog mere oplagt at videre specifiserer  
    
    \newpage
    
    % First task
    A. Brugeren, Instalatør, alle sensorene og DMI.
En bruger kan altid se indsamlet data, fra enten Sensorene i produktet eller DMI.

En Instalatør vil blive spurgt for en administrator kode inden vedkomende vil få adgang til administrerende settings, så som kalibrering eller data manipulation.

En sensor kan tilgå systemet og tilføje sensor data.

Produktet henter og opdatere data fra DMI hver 15 miniut.

B. Bruger og administrator.


%------------------------------------
\section{Usecase beskrivelse}
\begin{figure}
    \centering
    \includegraphics{}
    \caption{Caption}
    \label{fig:my_label}
\end{figure}
Systemet indeholder som usecase diagrammet anviser fire forskellige aktører:\\
En bruger, der kan få vist den indsamlede data, gennem en af de to mulige grænseflader, for henholdsvis den samlede Sensor data, eller den lokale udsigt opdateret fra DMI.\\
En installatør, der gennem en netbaseret forbindelse kan loggein i systemet, for efterfølgende at administrer produktet, ved f.eks. opsætning eller omkonfigurering, sensor kalibrering, eller datamanipulation.\\
En rakke sensorer der bidrager med forskellige målte data til systemets lokale information, som senere kan fremvises til brugeren.\\
Og en tilgang til DMI's kommende vejrudsigt, der bliver opdateret for brugerens oplevelse.
\\\\
\noindent
Produktets tænkte leveprocess tænkes som:\\
Brugeren køber produktet fra en producent eller distributør.\\
Installerer produktet i sit hjem ved hjælp af den medfølgende manual, og op-kobler produktets hovedmodul til en lokal netværksruter.\\
Efterfølgende kan brugeren tilgå hovedmodulet gennem dens uddelte lokale ipaddresse, fra f.eks. en lokal computer.\\
Igennem denne forbindelse til hovedmodulets, vil brugeren blive guidet gennem en førstegangs opsætnings wizzard.\\
Her efter kan denne forbindelse også bestyres eller konfigurere gennem denne samme tilgang.
    \newpage
    
\end{document}